% ECONOMICS_PROBLEM: {"id": "F13-230", "title": "Costs of geopolitical fragmentation", "jel_code": "F13", "source_id": "OP-0344"}
% !TeX program = xelatex
\documentclass[11pt,a4paper]{article}
\usepackage[margin=23mm]{geometry}
\usepackage{fontspec}
\setmainfont{Latin Modern Roman}
\usepackage{amsmath,amssymb,mathtools}
\usepackage{xcolor,enumitem,booktabs,longtable,array,xurl,hyperref}
\definecolor{muted}{HTML}{53636C}
\hypersetup{colorlinks=true,urlcolor=blue,linkcolor=blue}
\setlength{\parindent}{0pt}
\setlength{\parskip}{5pt}
\newcommand{\R}{\mathbb R}
\newcommand{\E}{\mathbb E}
\newcommand{\N}{\mathbb N}
\newcommand{\ind}{\mathbf 1}
\newcommand{\Var}{\operatorname{Var}}
\newcommand{\Cov}{\operatorname{Cov}}
\newcommand{\supp}{\operatorname{supp}}
\DeclareMathOperator*{\argmin}{arg\,min}
\DeclareMathOperator*{\argmax}{arg\,max}
\newcommand{\entrymeta}[3]{{\sffamily\small\color{muted}\textbf{\texttt{#1}}\quad|\quad #2\par #3\par}}
\newcommand{\Problemref}[1]{\texttt{#1}}
\newcommand{\JELref}[2]{\texttt{#2} (JEL \texttt{#1})}
\begin{document}
\section*{Costs of geopolitical fragmentation}
\textbf{Problem ID:} \texttt{F13-230}\quad \textbf{JEL:} \texttt{F13}\par
% BEGIN ECONOMICS STATEMENT
\label{problem:OP-0344}
\label{audited:ECON-078}
{\sffamily\small\color{muted}Original subject group: Trade, globalization and geoeconomics\par}
\label{definitions:problem:30007626}
\textbf{Audit classification:} Published research agenda. The official fragmentation research priority is verified. The detailed dated scenario and consumption-equivalent welfare problem are editorial; existence/uniqueness under stated monotonicity are routine mathematical facts.
\subsection*{Definitions and precise research target}
Fix a multi-country, multi-sector economy with trade and investment networks. A fragmentation scenario \(a\) specifies bilateral restrictions by partner, sector, asset type and implementation date. Firms may switch suppliers, enter, exit and reorganize production, subject to explicit adjustment technologies; knowledge diffusion and investment are included only through declared equations. Let \(C_{ct}(a)\) denote representative-country or distribution-weighted consumption in equilibrium. Define the discounted welfare loss using a consumption-equivalent factor \(\lambda_c\):
\[E\sum_{t\geq0}\beta^t u_c((1-\lambda_c)C_{ct}(0))=E\sum_{t\geq0}\beta^t u_c(C_{ct}(a)),\]
where \(\beta\in(0,1)\), strictly increasing utility \(u_c\) and the no-fragmentation path \(0\) are fixed, and sums converge. Require \(\lambda_c<1\) and assume the comparison utility lies in the range generated by proportional baseline consumption; negative values indicate a welfare gain. The task is to identify or bound \(\lambda_c\) and world productivity effects under plausible adaptations and trade rerouting. State restrictions identifying substitution elasticities, fixed reorganization costs and geopolitical selection separately. Report sensitivity to bloc definitions, connector countries and the probability of future policy reversal. Distinguish already observed partial fragmentation from hypothetical complete decoupling. Model-based scenario estimates answer conditional counterfactuals; they are not identified causal global losses merely because a calibrated model reproduces baseline trade shares.

\textbf{Additional formal definitions and scope (editorial).} A fragmentation scenario is a dated matrix of tariffs, quantitative limits, investment eligibility and knowledge-flow restrictions; these cannot be represented by a generic bloc label alone. Use per capita consumption or a type-weighted consumption path consistently. Define \(J_c(\lambda)=E\sum_t\beta^t u_c((1-\lambda)C_{ct}(0))\) on \(\lambda<1\), with positive baseline consumption and finite sums. If \(J_c\) is continuous strictly decreasing and scenario utility belongs to its range, the consumption-equivalent loss exists and is unique; this is a definition/result under assumptions, not an open existence conjecture. Heterogeneous households require separate equivalents or an explicitly cardinal social welfare function. Trade diversion, rerouting and adjustment can change both equilibrium prices and quantities; a comparative-static scenario and a causal estimate of actual policy are different estimands.

For the identification part, declare an admissible class \(\Theta\) of structural models, including preferences, technologies, shock laws, measurement/sampling rules and any equilibrium selector. If \(O\) denotes the complete observable history and \(T(\theta)\) the target defined above, the identified set is \(\mathcal I_T=\{T(\theta):\theta\in\Theta,\ \mathcal L_\theta(O)=\mathcal L_{\rm obs}(O)\}\); point identification means this set is a singleton. A \(do(a)\) comparison replaces stated policy/technology equations while fixing the declared exogenous law and re-solving the remaining equations. These definitions do not assert that an identification theorem holds without further restrictions.
\subsection*{Variants and assumption changes}
\begin{enumerate}
\item \textbf{Static specified blocs (editorial).} Impose a dated trade-cost matrix and fixed technologies; compute conditional country consumption-equivalent losses.
\item \textbf{Dynamic adaptation and reversal (editorial).} Allow supplier switching, knowledge diffusion and uncertain policy reversal; identify or bound loss paths and robustness across connector-country arrangements.
\end{enumerate}
\subsection*{References and source audit}
\begin{enumerate}
\item Official January 9, 2026 web version; locator: Interconnected World / 2026 fragmentation priority. Broad agenda; exact model editorial. URL: \url{https://www.bankofengland.co.uk/research/bank-of-england-agenda-for-research}.
\end{enumerate}
\begin{enumerate}\raggedright
\item\hypertarget{definitions:source:30007626:1}{} Bank of England. 2026. \emph{Bank of England Agenda for Research, 2025–2028}. Interconnected World / 2026 fragmentation priority.
\url{https://www.bankofengland.co.uk/research/bank-of-england-agenda-for-research}.
\end{enumerate}

\subsection*{Integrated refinement: ECON-078}
{\sffamily\small\color{muted}Welfare effects of trade fragmentation and supply-chain reorientation\par}
This refinement adopts the additional source conventions
(page~\pageref{audited:conventions}), subject to its local restrictions.
\textbf{Money-metric distribution of fragmentation and sourcing changes.}
For a dated bilateral trade-cost change and a selected equilibrium, also report
the baseline-price equivalent variation
$EV_r(z)=e_r(p_0,u_r(z))-e_r(p_0,u_r(0))$ where the expenditure function is
well defined. Compare its distribution with the existing proportional
consumption-equivalent losses; the two are different normalization conventions.
Fix welfare populations, country/household weights and baseline resources.
Permit sourcing, entry, location, third-country diversion and adjustment costs
to respond rather than hold sourcing shares fixed. Identification or sensitivity
bounds must disclose trade elasticities, institutional restrictions and the
equilibrium selector; a calibrated scenario remains conditional on these inputs.
\subsection*{Supplied source audit for ECON-078}
Local [S1], [S2], etc. in this audit and the references immediately below
belong to ECON-078, independently of the original entry's reference numbering.
[S1] and [S2] support fragmentation, supply chains and developing-country trade effects. The scenario-specific equivalent-variation calculation is editorial.
\subsection*{References for integrated refinement ECON-078}
{\footnotesize\raggedright
\par\noindent\textbf{[S1]} David Atkin, Amit Khandelwal, Laura Boudreau, Rafael Dix-Carneiro, Isabela Manelici, Pamela Medina, Brian McCaig, Ameet Morjaria, Luigi Pascali, Heitor Pellegrina, Bob Rijkers, and Meredith Startz (2025). \emph{International Trade}. VoxDevLit 4(2), February 2025.\par \textit{Verified locator / source role:} Primary publisher page checked for review identity, authors, version and background. The specific published agenda is corroborated in [S2]; the exact formal target is editorial.\par \url{https://voxdev.org/voxdevlit/international-trade}\par\smallskip
\par\noindent\textbf{[S2]} Oliver Hanney / VoxDev (living compilation). \emph{Evidence gaps: Key unanswered questions in development economics}. Accessed 8 October 2026. The page currently covers 20 topics; the ``16-topics URL is a historical slug.\par \textit{Verified locator / source role:} Topic heading: International Trade. Supports the broad research agenda; this is not a verbatim quotation or an attribution of the displayed mathematical target.\par \url{https://voxdev.org/topic/evidence-gaps-key-unanswered-questions-across-16-topics-development-economics}\par\smallskip
}

\subsection*{Common conventions: Source mathematical conventions}

These conventions apply to each entry unless explicitly replaced.

\textbf{Models and identification.} A primitive model \(m\) specifies all probability laws, feasible choices, information, equilibrium and measurement rules used by its target. The admissible class \(\mathcal M\) is fixed independently of outcomes. If \(P_O(m)\) is its observable probability law and \(\tau(m)\) the target, its identified set at observable law \(P\) is
\[
 \mathcal I_\tau(P)=\{\tau(m):m\in\mathcal M,\ P_O(m)=P\}.
\]
Point identification means that this set is a singleton. An empty set signals incompatibility of data and assumptions. Coordinate bounds need not describe the joint set. Bounds are sharp only if every retained target is attainable or arbitrarily approachable by compatible models. Finite bounds require explicit restrictions on unobserved quantities; unrestricted confounding, hidden wealth, extrapolation or arbitrary equilibrium selection can yield uninformative sets.

\textbf{Causal interventions.} A potential outcome \(Y_i(p)\) is the outcome under a complete policy regime \(p\), including timing, financing, information and market spillovers. The target population and units are fixed before treatment. With interference, use the complete assignment vector or a justified exposure mapping; individual no-interference notation alone cannot identify an economy-wide intervention. A difference of mean potential outcomes does not require the cross-world joint distribution. Conditioning on post-treatment migration, survival, enrollment or employment changes the estimand and needs separate justification.

\textbf{Statistical statements.} An estimator, confidence set, forecast or test specifies a sampling law, model class and loss/coverage criterion. Uniform \((1-\alpha)\)-coverage means \(\inf_{m\in\mathcal M}\Pr_m\{\tau(m)\in C_n\}\geq1-\alpha\), or an explicitly asymptotic version. The whole parameter space trivially has coverage, so informativeness requires a stated width, risk or power objective. A forecasting improvement is relative to a named benchmark, information set and evaluation population.

\textbf{Equilibrium and optimization.} An equilibrium comprises feasible plans, optimality, beliefs where needed, budget constraints and all market/resource clearing. The primitives or a stated benchmark define each correspondence \(\mathcal E(p;m)\). Nonemptiness is assumed or proved, never inferred just from compact policy sets. A selected-equilibrium target uses a declared selection rule; a robust target quantifies over all permitted equilibria. Use suprema or infima unless attainment is established. An \(\varepsilon\)-optimal policy is feasible and within \(\varepsilon\) of the finite optimum value. Report an empty feasible set separately.

\textbf{Welfare and accounting.} Normative utility, interpersonal normalization, social weights, numeraire, discounting and the use of public receipts are primitives. Behavioral utility need not coincide with the welfare criterion. Household welfare includes ownership income and rents once; adding profits again to compensating variation generally double-counts them. A monetary equivalent is defined by an explicit transfer and price convention, with monotonicity and range conditions. Average effects, mechanism contributions, transfers and welfare losses are different targets. Infinite sums require convergence and dynamic feasible plans require terminal or no-Ponzi conditions.

\textbf{Source versions.} A working-paper date, revision date and journal publication year are distinguished. A DOI for a journal version is not automatically the DOI of an earlier manuscript. Locators refer to the stated version and printed pages where specified. A metadata-only check does not verify a claim about a full-text conclusion. Every entry's source subsection narrows the provenance claim accordingly.

\subsection*{Common conventions: Additional-source status and mathematical conventions}

\label{audited:conventions}

Of the 100 supplied ECON entries, 65 distinct problems appear here; 32 close
matches refine earlier entries, and three duplicate questions map to existing
entries. Appendix~\ref{app:audited-crosswalk} lists every destination.
The attachment's P/R/S/U distinctions and source audits are retained. Its
precise mathematical formalizations are editorial unless the individual audit
says otherwise. Candidate subsidy targets ECON-004--006 retain their U flags;
the resolved approval-core question ECON-007 updates OP-0202 above.
The following supplied conventions govern both this part and the attributed
ECON refinements elsewhere in the catalogue, subject to local restrictions.

\textbf{Parameterized questions.} A problem family lists inputs that must be fixed before an instance is evaluated: population, horizon, policies, information, data law, model class and welfare weights. These are required choices, not quantities the citations are assumed to specify. Undefined applications should not be treated as identified estimands.

\textbf{Indices and integrability.} Sets of agents, firms and regions are finite unless a population measure is explicitly used. \(i,j\) index units, \(r\) regions, \(t\) dates and \(h\) horizons. \(\mathbb E\) and \(\Pr\) refer to the declared population and law; \(\mathbf1\{A\}\) is an indicator. Every displayed expectation and discounted sum needs existence in the stated domain. Infinite horizons use a discount in \((0,1)\) and a convergence condition; finite horizons may use discount one.

\textbf{Optimization and existence.} A displayed \(\max\), \(\min\), or \(\arg\max\) is conditional on attainment. For finite feasible menus this is immediate; general spaces need continuity and compactness/coercivity or another existence argument. Without attainment, the target is the supremum/infimum and arbitrarily near-optimal feasible choices. \(\inf\varnothing=+\infty\). Empty feasibility and an unidentified target are different issues.

\textbf{Potential outcomes.} \(Y_i(z)\) is the outcome under a specified intervention, not the observed conditional mean among people choosing \(z\). In binary comparisons \(\Delta Y_i=Y_i(1)-Y_i(0)\). Specify assignment, treatment versions, compliance, information and observation horizon. Interference is allowed under a full policy vector; compressed exposure notation needs a justified mapping. No interference, consistency and overlap are substantive assumptions, not automatic consequences of notation.

\textbf{Populations and observation.} Fix baseline eligibility and population weights. Conditioning on post-policy employment, survival, relocation or program uptake can change the population being compared. Death is not ordinary loss to follow-up; outcomes truncated by death need a composite convention, joint survival target, or explicitly defined survivor stratum. Fertility-changing interventions need a population functional rather than an unexplained fixed descendant cohort. Measurement and missingness models must be supplied.

\textbf{Identification.} A structural model \(m\in\mathcal M\) implies observable law \(P_m\) and target \(\theta(m)\). Identification means
\[
 P_{m_1}=P_{m_2}\ \Longrightarrow\ \theta(m_1)=\theta(m_2)
 \quad\text{for all }m_1,m_2\in\mathcal M .
\]
Without identification, the target is
\[
 \Theta(P;\mathcal M)=\{\theta(m):m\in\mathcal M,\ P_m=P\}.
\]
Writing \(\theta(P)\) before establishing identification can hide incompatible latent models. Confidence sets additionally need a sampling law, confidence level, asymptotic sequence and uniformity class. Point identification, coverage and informative precision are separate properties.

\textbf{Mediation.} Total effects of randomized assignment are distinct from cross-world natural indirect effects. Belief/effort-channel effects require a meaningful mediator intervention and additional measurement, exclusion or mediation assumptions. Randomization of the initial policy alone does not supply those assumptions. Report bounds or interventional alternatives when the stronger target is unsupported.

\textbf{Equilibrium.} A counterfactual \(W(z)\), \(p(z)\), or \(\mu_t^z\) requires individual optimization, feasibility, government/resource budgets, market clearing and state-transition laws. State equilibrium existence and selection, or report ranges over compatible equilibria. Initial-state integration includes subsequent endogenous transitions; current-state integration must use the policy-dependent distribution. Reduced-form invariance after policy is not assumed.

\textbf{Welfare accounts.} Scores, utility, health, dollars and ecological quantities require explicit conversion before addition. Fees, taxes, subsidies, wages and extortion payments are transfers between parties; distributional weights can make them welfare relevant, but their face value is not automatically a resource loss. State ownership, geographic boundaries, foreign-income weights and fiscal finance. Do not add profits to household welfare already including dividends, or count land capitalization and the underlying benefit twice. A benefit-cost expression must distinguish gross benefits from net welfare and subtract every real cost exactly once.

\textbf{Derivatives and risk.} Log derivatives require positive arguments; local responses require admissible perturbations and specified equilibrium branches. Differentiation under expectations needs regularity/domination, and boundaries may allow only one-sided derivatives. Risk uses a specified law; ambiguity uses a declared nonempty law set on a common history space. Adaptive policies must be nonanticipative. A robust criterion is an editorial choice unless attributed explicitly.

\textbf{Scope overlaps.} Allocation, collective choice, contracts, household bargaining and coordinated policy overlap economics with optimization and game theory. They are retained because they appeared in the original catalogue; no claim of a strictly game-theory-free selection is made.
% END ECONOMICS STATEMENT
\end{document}
